
          %%%%% ~~~~~~~~~~~~~~~~~~~~ %%%%%

\chapter{Spin, magnetic dipole moment, the Stern-Gerlach experiment, and
the Pauli-Schr\"{o}dinger equation as a non-relativistic limit of the Dirac equation}
\setcounter{theorem}{0}
\setcounter{equation}{0}

%\cite{vanDerVaart1996}
%\cite{Kosorok2008}

%\renewcommand{\theenumi}{\alph{enumi}}
%\renewcommand{\labelenumi}{\textnormal{(\theenumi)}$\;\;$}
\renewcommand{\theenumi}{\roman{enumi}}
\renewcommand{\labelenumi}{\textnormal{(\theenumi)}$\;\;$}

          %%%%% ~~~~~~~~~~~~~~~~~~~~ %%%%%

In this section, we discuss the following:
\begin{itemize}
\item
    how the Pauli-Schrödinger equation arises through canonical quantization
    of the Hamiltonian of an electrically charged particle,
    possibly with non-zero intrinsic spin, inside an electromagnetic field, 
\item
    how the Pauli-Schrödinger equation can be used to explain the observations
    of the Stern-Gerlach experiment,
\item
    how the Pauli-Schrödinger equation can be regarded as a non-relativistic limit
    of the Dirac equation, and
\item
    the expression of the magnetic dipole moment of a particle
    in terms of the intrinsic spin of the particle.
\end{itemize}

          %%%%% ~~~~~~~~~~~~~~~~~~~~ %%%%%

\vskip 0.5cm
\section{Spin}

          %%%%% ~~~~~~~~~~~~~~~~~~~~ %%%%%

          %%%%% ~~~~~~~~~~~~~~~~~~~~ %%%%%

          %%%%% ~~~~~~~~~~~~~~~~~~~~ %%%%%

          %%%%% ~~~~~~~~~~~~~~~~~~~~ %%%%%

